\scanpage{50}
On a
\[
\begin{aligned}
(16)\quad&
\left\{\begin{aligned}
\text{a)}\quad&\text{Si }\alpha:\deleted{\unsure{\alpha}}\Delta_{r'}\longrightarrow\Delta_r\text{ est tel que }\alpha(0)=0,\text{ alors}\\
&X_{\Delta_r}\xrightarrow{\alpha^*}X_{\Delta_{r'}}\text{ est un morphisme de rev.}\\
&\text{étale (en particulier}\\
&\left\{\begin{aligned}
X_{\Delta_r}&\longrightarrow X_{\Delta_0}\\
(x_0,\ldots,x_r)&\longmapsto x_0\quad\text{ou }(x,i_0,\ldots,i_r)\longmapsto x
\end{aligned}\right.
\end{aligned}\right.\\
&\qquad\text{est un morphisme de rev. étale}
\end{aligned}
\]
En fait, la \deleted{fibre de} restriction de $X_{\Delta_r}$ au
dessus des composantes $X_{i'_*}$ \inserted{$=X_{i'_0}\times\{i'_*\}$} de $X_{\Delta_{r'}}$
(où $i'_* = \{i'_0\leq\cdots\leq i'_{r'}\}$) \deleted{est formée des} s'identifie
au rev étale constant de fibre type l'ens
des drapeaux \inserted{$\unsure{i_*|d_*}$} de $I$, de type $\Delta_r$, qui
s'envoient sur $i'_*$ par $\alpha^*$ :
\[
(17)\quad\alpha^{*^{-1}}\left(\underbrace{X_{i'_*}}_{\substack{\Vert\\X_{i'_0}\times\{i'_*\}}}\right)
\simeq X_{i'_*}\times\{i^1_*\in I_{\Delta_r}\mid\alpha^*(i_*)=i'{}^*\}
\]
Quand aux opérations semi-simpliciales
générales (13) (sans condition $\alpha(\unsure{\omega|0})=0$ sur
$\alpha$) on peut dire seulement qu'elles sont
des \uline{imm. locales}. En fait, tous les
$X_{\Delta_r}$ sont au dessus de $X$, \deleted{donc}
\[
(18)\qquad X_{\Delta_r}\longrightarrow X\qquad\text{\uline{imm. \unsure{focale|locale}}}.
\]
\margin{49}
